\documentclass{article}

\usepackage[printqrbox=false,printhint=true,printanswer=true,printmarkingguide=false]{unswalgos}

\usepackage{tikz}
\usetikzlibrary{patterns}
\usetikzlibrary{shapes,fit}
\usepackage{tkz-fct}
\usepackage{wrapfig}
\usepackage{subfig}

\usepackage{mathtools}
\usepackage{amssymb}
\usepackage{booktabs,multicol,multirow}
\usepackage{wasysym}
\usepackage{tcolorbox}

\DeclareMathOperator*{\argmax}{arg\,max}
\DeclareMathOperator*{\argmin}{arg\,min}
\DeclareMathOperator{\NAND}{NAND}
\DeclareMathOperator{\AND}{AND}
\DeclareMathOperator{\OR}{OR}
\DeclareMathOperator{\NOT}{NOT}

\usepackage{xspace}

\fancyfoot[L]{\leftmark}
\fancyfoot[R]{\rightmark}

% This enables new paragraphs without indentation
\usepackage[parfill]{parskip}

\newcommand{\sem}{23T2}
\newcommand{\semester}{Term 2, 2023}
\SubjectNo{COMP3121/9101}
\newcommand{\taskname}{Assignment 2, Question 3}
\Institution{Algo Rhythm, z1234567} % Replace this with your name and zID


\begin{document}

\setcounter{question}{2}

\begin{Question}

\begin{Subquestion}
\textbf{[10 marks]} 
Santa's assistant has been laid off and it is now your job to ensure that as many \textbf{wrapped} presents get delivered to shopping centres as possible.

Santa has an infinite supply of \textbf{unwrapped} presents at the North Pole, and can distribute them to $k$ workshops across the globe. At each workshop $i$, a maximum of $a_i$ presents can be wrapped, and the North Pole can distribute to it at most $d_i$ presents.

The workshops wrap the presents, and then deliver these to $m$ shopping centres. Each shopping centre $j$ can hold up to $b_j$ packages, however, due to some legacy contractual agreements, each shopping centre can only receive presents from certain workshops. You are given the contractual agreements as a two dimensional array $C$, where $C[i][j]$ denotes that workshop $i$ has an agreement with shopping centre $j$.

Design an algorithm that runs in $O((k+m)k^{2}m^{2})$ time which finds the maximum number of wrapped presents that can be delivered to shopping centres.
\end{Subquestion}

\begin{answer}
Create a flow network as below:
\begin{itemize}
    \item A source $s$ and a sink $t$.
    \item A vertex $w_i$ for the $i$ th workshop.
    \item A vertex $sh_i$ for the $i$ th shopping center.
    \item An edge from $s$ to the each $w_i$ of capacity $\min{a_i, d_i}$.
    \item For every pair of workshops and shopping centers, if the $i$ th workshop and the $j$ th shopping center has an agreement, link an edge from $w_i$ to $sh_j$ with infinity capacity.
    \item An edge from each $sh_i$ to $t$ with capacity  $b_j$.
\end{itemize}
Run the Edmonds-Karp algorithm from to find the maximum flow $f$ from $s$ to $t$. $f$'s value is the number of wrapped presents that can be delivered to shopping centers.

The algorithm is correct since our flow network can well represent the delivery process of the presents and the constraints. the source $s$ is the north pole, 
the edge from the north pole to the workshops is with the capacity $\min{a_i, d_i}$, represent that we can only distribute to the workshop $d_i$ presents and the workshop can wrap at most $a_i$ of them.
The edge from the workshops to the shopping centers follows the agreements, and from the shopping center to the sink, the capacity is to ensure that only $b_i$ present can be stored in the shopping center.
Thus the maximum flow from $s$ to $t$ is the maximum present that can be delivered to shopping centers.

In the flow network, we have $2 + k+m = O(k+m)$ vertices and at most $k\times m + k + m =O(km)$ edges. Thus the runtime of the Edmonds-Karp Algorithm is $O(VE^2) = O((k+m)k^2m^2)$.  
\end{answer}

\begin{Subquestion}
\textbf{[6 marks]} 
On Christmas Eve, you realise that you should ensure that each of the $n$ children on Santa's nice list receive at least one present. Children are willing to travel at most $D$ distance to a shopping centre to receive a present. You are given the location of all children and shopping centres.

Design an algorithm that runs in $O(nm(k+n))$ time which finds the maximum number of children on the nice list who can receive at least one present. 
\end{Subquestion}

\begin{answer}
Create a flow network as below:
\begin{itemize}
    \item A source $s$ and a sink $t$.
    \item A vertex $w_i$ for the $i$ th workshop.
    \item Two vertex $sh1_i$ and $sh2_i$ for the $i$ th shopping center.
    \item A vertex $c_i$ for the $i$ the child.
    \item An edge from $s$ to the each $w_i$ of capacity $\min{a_i, d_i}$.
    \item For every pair of workshops and shopping centers, if the $i$ th workshop and the $j$ th shopping center has an agreement, link an edge from $w_i$ to $sh1_j$ with infinity capacity.
    \item An edge from $sh1_j$ to $sh2_j$ with capacity $b_j$.
    \item For ever child-shopping center pair, if the distance between child $i$ and shopping center $j$ is less than $D$, link $sh2_j$ and $c_i$ with an edge of capacity $1$.
    \item An edge from each $c_i$ to $t$ with capacity  $1$.
\end{itemize}
Run the Ford-Fulkson algorithm from to find the maximum flow $f$ from $s$ to $t$. $f$'s value is the number of child who can receive at least 1 present.

The algorithm is correct since our flow network can well represent the delivery process of the presents and the constraints. The network is almost the same as the one 
in 3.1, but we use the edge between $sh1_i$ and $sh2_i$ to represent the number of presents can be stored in the shopping center.
And we use the edges between children with shopping centers to represent the purchase relations.
Thus the maximum flow from $s$ to $t$ is the maximum number of children who can receive presents.

In the flow network, we have at most $k\times m + k + m + n + m \times n =O(m(k+n))$ edges. And the maximum value of the max flow is $n$, thus the runtime of the Edmonds-Karp Algorithm is $O(\mid E\mid f) = O(nm(k+n))$.  

\end{answer}

\begin{Subquestion}
\textbf{[4 marks]} 
Let the farthest distance from any child to a shopping centre be $F$.

Design an algorithm that runs in $O(nm(k+n)\log{F})$ time which finds the minimum value of $D$ such that all $n$ children on the nice list receive at least one present, or returns that no such $D$ is possible.

\begin{answer}
We use a binary search to solve the problem here. Let the minimum value of $D$ be $D_{min}$, $left = 0, right = F$. We do binary search in the range $[0,F]$ to search for $D_{min}$. 
Every time we search a value $mid$, we let $D=mid$ and use 3.2's algorithm to find the maximum number of children who 
can receive a present. If the number equal to $n$, we let $right = mid$, else we let $left = mid$. When $right - left \le 1$, the value of right 
is $D_{min}$ we need to find.

The algorithm is true since it's clear that if a $D$ can satisfy the problem's requirement, any other $D' >D$ can satisfy too, and if $D$ is unsatisfiable, any other $D' < D$ is also unsatisfiable. Thus we can use binary search to search $D_{min}$.
And in 3.2 we have already proved our algorithm to check whether a $D$ is satisfiable is correct, thus the whole algorithm can give the correct answer.

The time complexity is $O(nm(k+n)\log F)$ since each check algorithm runs in $O(nm(k+n))$ time, and in binary search, we need to do the check for $O(\log F)$ times.
\end{answer}

\end{Subquestion}

\end{Question}

\end{document}